\section*{第 \ref{ch:b2:plant-life-cycles} 章——陆生植物的生活史与生殖}

\begin{solution}{exo:b2:plant-life-cycles:1}
\omterm{def:b2:plant-life-cycles:cycles}{配子体}：\omterm{def:b2:meiosis-heredity:meiosis}{单倍体}的多细胞世代，由孢子经有丝分裂长成，
通过有丝分裂产生配子。\omterm{def:b2:plant-life-cycles:cycles}{孢子体}：\omterm{def:b2:meiosis-heredity:meiosis}{二倍体}世代，
由受精卵经有丝分裂长成，通过\omterm{def:b2:meiosis-heredity:meiosis}{减数分裂}产生孢子。
孢子：由\omterm{def:b2:meiosis-heredity:meiosis}{减数分裂}产生、无需融合即可生长的\omterm{def:b2:meiosis-heredity:meiosis}{单倍体}细胞。
配子：（在植物中）由有丝分裂产生的\omterm{def:b2:meiosis-heredity:meiosis}{单倍体}细胞，
必须与另一个融合。
\end{solution}

\begin{solution}{exo:b2:plant-life-cycles:2}
叶细胞为 $2n = 1260$：孢子 630，\omterm{def:b2:plant-life-cycles:fern}{原叶体}细胞 630，精子
630，卵 630，受精卵 1260。
\end{solution}

\begin{solution}{exo:b2:plant-life-cycles:3}
绿色垫状体：\omterm{def:b2:plant-life-cycles:cycles}{配子体}。蒴柄与蒴：\omterm{def:b2:plant-life-cycles:cycles}{孢子体}。蕨
叶：\omterm{def:b2:plant-life-cycles:cycles}{孢子体}。\omterm{def:b2:plant-life-cycles:fern}{原叶体}：\omterm{def:b2:plant-life-cycles:cycles}{配子体}。\omterm{def:b2:plant-life-cycles:heterospory}{花粉}粒：雄性
\omterm{def:b2:plant-life-cycles:cycles}{配子体}。种子贮藏物（松树中）：雌性\omterm{def:b2:plant-life-cycles:cycles}{配子体}。
种子的胚：下一代\omterm{def:b2:plant-life-cycles:cycles}{孢子体}。
\end{solution}

\begin{solution}{exo:b2:plant-life-cycles:4}
\omterm{def:b2:meiosis-heredity:meiosis}{单倍体}型\omterm{def:b2:plant-life-cycles:cycles}{生活史}（\emph{Chlamydomonas}）：受精卵立即进行\omterm{def:b2:meiosis-heredity:meiosis}{减数分裂}。
\omterm{def:b2:meiosis-heredity:meiosis}{二倍体}型\omterm{def:b2:plant-life-cycles:cycles}{生活史}（动物、\emph{Fucus}）：\omterm{def:b2:meiosis-heredity:meiosis}{减数分裂}产生配子。
单倍–\omterm{def:b2:meiosis-heredity:meiosis}{二倍体}型\omterm{def:b2:plant-life-cycles:cycles}{生活史}（\omterm{def:b2:plant-life-cycles:moss}{苔藓}、\omterm{def:b2:plant-life-cycles:fern}{蕨类}、种子植物和许多藻类）：\omterm{def:b2:meiosis-heredity:meiosis}{减数分裂}
在\omterm{def:b2:plant-life-cycles:cycles}{孢子体}的\omterm{def:b2:plant-life-cycles:fern}{孢子囊}中产生孢子。
\end{solution}

\begin{solution}{exo:b2:plant-life-cycles:5}
$0.03/10^{-4} = \qty{300}{s}$，即五分钟。在 \qty{20}{min}
内可移动 \qty{12}{cm}。受精只在相距几厘米的植株间进行，
因此\omterm{def:b2:plant-life-cycles:moss}{苔藓}长成致密的垫状群落，使雌雄
（或两种器官）都处于可及范围之内。
\end{solution}

\begin{solution}{exo:b2:plant-life-cycles:6}
$v_s = 2\times(1.5\times 10^{-5})^2\times 1099\times 9.81/(9\times
1.8\times 10^{-5}) = \qty{0.030}{m/s}$，\qty{3}{cm/s}；距离 $uh/v_s
= 2\times 0.5/0.03 = \qty{33}{m}$。种子：$2\times 20/0.8 =
\qty{50}{m}$——较重而带翅的传播单位从高处脱落后
传播的距离大致相同。
\end{solution}

\begin{solution}{exo:b2:plant-life-cycles:7}
$200\times 40\times 64 = \num{512000}$ 个孢子（每叶）；$5.1$
个\omterm{def:b2:plant-life-cycles:fern}{原叶体}；$0.1$ 个幼\omterm{def:b2:plant-life-cycles:cycles}{孢子体}（每叶）。
\end{solution}

\begin{solution}{exo:b2:plant-life-cycles:8}
种子即是被保留的\omterm{def:b2:plant-life-cycles:heterospory}{胚珠}：雌\omterm{def:b2:plant-life-cycles:cycles}{配子体}必须留在
母体上，受包被保护并获得营养，而雄\omterm{def:b2:plant-life-cycles:cycles}{配子体}必须
迁移到它那里。这需要两类孢子，各有其去向。同型孢子植物
的唯一孢子必须被释放，才能长成自由的
两性\omterm{def:b2:plant-life-cycles:cycles}{配子体}；若保留它，也就保留了雄性功能，
并迫使每株植物自体受精，于是
便没有需要迁移的东西了。
\end{solution}

\begin{solution}{exo:b2:plant-life-cycles:9}
孢子：$\tfrac43\pi(1.5\times 10^{-5})^3\times 1100 =
\qty{1.6e-11}{kg}$，\qty{16}{ng}；能量 $0.5\times 1.6\times
10^{-8}\times 2\times 10^{4} = \qty{1.6e-4}{J}$。种子：\qty{5}{mg}，
$0.5\times 5\times 10^{-3}\times 2\times 10^{4} = \qty{50}{J}$ ---
是其 $3\times 10^{5}$ 倍。孢子只能长出一个小型
光合鳞片，见光前别无他物；种子则可在黑暗中长出
根与芽，能等待一个季节，并能在
被吃掉或干燥之后存活。
\end{solution}

\begin{solution}{exo:b2:plant-life-cycles:10}
$10^{11}/(10^{6}\times 20) = \num{5000}$ \omterm{def:b2:plant-life-cycles:heterospory}{花粉}粒每立方米；通量
$5000\times 10^{-7}\times 2 = \qty{1e-3}{}$ \omterm{def:b2:plant-life-cycles:heterospory}{花粉}粒每秒；约
\qty{1000}{s}，即四分之一小时，每粒。\omterm{prop:b2:plant-life-cycles:conifer}{球果}必须开放
并保持可受粉状态数小时至数日，恰与雄\omterm{prop:b2:plant-life-cycles:conifer}{球果}
散粉同步——传粉时间精确到周。
\end{solution}

\begin{solution}{exo:b2:plant-life-cycles:11}
二倍性掩盖了\omterm{def:b2:meiosis-heredity:vocabulary}{隐性}有害\omterm{def:b2:genome-diversification:mutation}{突变}；具有维管组织的\omterm{def:b2:meiosis-heredity:meiosis}{二倍体}
植株可长得很大且寿命很长，高度既有利于光捕获，也有利于
孢子传播；\omterm{def:b2:plant-life-cycles:cycles}{孢子体}高处产生的孢子随空气扩散，
而配子则必须游动。\omterm{def:b2:meiosis-heredity:meiosis}{单倍体}世代
之所以持续存在，是因为\omterm{def:b2:meiosis-heredity:meiosis}{减数分裂}与受精都需要
\omterm{def:b2:meiosis-heredity:meiosis}{单倍体}阶段；\omterm{def:b2:plant-life-cycles:heterospory}{花粉}粒与胚囊是最简结构，
用以将配子彼此输送并为胚胎提供营养。
\end{solution}

\begin{solution}{exo:b2:plant-life-cycles:12}
他通过培养孢子并在显微镜下跟踪观察，得以确定
孢子长成一个带性器官的小型
植体；胚由受精卵产生，该受精卵位于
\omterm{def:b2:plant-life-cycles:moss}{颈卵器}内，\omterm{def:b2:plant-life-cycles:cycles}{孢子体}由此长成，并且
同样的器官以退化形式存在于裸子植物\omterm{def:b2:plant-life-cycles:heterospory}{胚珠}中：
\omterm{def:b2:plant-life-cycles:cycles}{世代交替}表现为植体与器官的序列。他无法
知道在细胞水平上是什么区分了两个世代。
染色体计数表明，两种植体的染色体数目相差
两倍，\omterm{def:b2:meiosis-heredity:meiosis}{减数分裂}发生于孢子形成，而受精发生于
卵细胞处，由此解释了\omterm{def:b2:plant-life-cycles:cycles}{世代交替}
为何必然发生。
\end{solution}

\begin{solution}{pb:b2:plant-life-cycles:1}
\textbf{1.} $10\times 200\times 40\times 64 = 5.1\times 10^{6}$
孢子。
\textbf{2.} $\tfrac43\pi(1.5\times 10^{-5})^3\times 1100 =
\qty{1.6e-11}{kg}$ 每个孢子；总量 $8\times 10^{-5}$ 千克，\qty{80}{mg}。
\textbf{3.} $v_s = 2\times 2.25\times 10^{-10}\times 1099\times
9.81/(1.62\times 10^{-4}) = \qty{0.030}{m/s}$.
\textbf{4.} $x = 1.5\times 0.5/0.030 = \qty{25}{m}$；圆盘
面积 $\pi\times 25^2 = \qty{2000}{m^2}$；约
\num{2600} 个孢子每平方米。
\textbf{5.} 每秒数厘米的上升气流超过 $v_s$，因此
被湍流卷起的孢子可在空中停留数小时，
并可传播数百千米；若其\omterm{def:b2:plant-life-cycles:fern}{原叶体}能自体受精
（它同时具有两种器官），单个孢子即可建立种群，
这正是偏远岛屿\omterm{def:b2:plant-life-cycles:fern}{蕨类}区系丰富的原因。
\textbf{6.} 孢子：$0.5\times 1.6\times 10^{-8}\times 2\times 10^{4} =
\qty{1.6e-4}{J}$；总量：\qty{820}{J}；一粒松种子：\qty{50}{J}——
整个孢子总量相当于十六粒种子。
\textbf{7.} $5.1\times 10^{6}/2\times 10^{4} = 256$ 个\omterm{def:b2:plant-life-cycles:fern}{原叶体}。
\textbf{8.} $1 - 0.8^{6} = 1 - 0.26 = 0.74$.
\textbf{9.} $256\times 0.74\times 0.5\times 0.02 = 1.9$ 株成年\omterm{def:b2:plant-life-cycles:fern}{蕨类}
每个季节。
\textbf{10.} $30\times 1.9 \approx 57$：远多于稳定种群的一个
替换量，因此存活率数据
偏乐观——在饱和的栖息地中，多数幼\omterm{def:b2:plant-life-cycles:cycles}{孢子体}死于
拥挤与遮荫。
\textbf{11.} $10^{-3}/10^{-4} = \qty{10}{s}$。许多\omterm{def:b2:plant-life-cycles:fern}{蕨类}使
\omterm{def:b2:plant-life-cycles:moss}{精子器}与\omterm{def:b2:plant-life-cycles:moss}{颈卵器}成熟时间错开，或释放一种激素
（\omterm{def:b2:plant-life-cycles:moss}{精子器}诱导素），使邻近的幼\omterm{def:b2:plant-life-cycles:fern}{原叶体}发育为雄性，
从而使精子来自另一个个体。
\textbf{12.} 它只有一个细胞厚，无根、无保护，需要
液态水完成受精，仅存活数周：几乎全部
\omterm{def:b2:plant-life-cycles:cycles}{生活史}的死亡都集中在它身上（每 $2\times 10^{4}$ 个孢子中
仅一个能成活；其中四分之一从未遇到湿润的夜晚）。
\textbf{13.} $v_s = 2\times 6.25\times 10^{-10}\times 599\times
9.81/(1.62\times 10^{-4}) = \qty{0.045}{m/s}$.
\textbf{14.} $3\times 15/0.045 = \qty{1000}{m}$.
\textbf{15.} $10^{11}/(2\times 10^{7}) = \num{5000}$ 每立方米。
\textbf{16.} $5000\times 10^{-7}\times 2 = \qty{1e-3}{s^{-1}}$：每小时
3.6 粒；首粒约在 \qty{1000}{s} 后到达，即四分之一
小时。
\textbf{17.} 顺风五千米处，\omterm{def:b2:plant-life-cycles:heterospory}{花粉}云已向两侧和
上方扩散，大部分\omterm{def:b2:plant-life-cycles:heterospory}{花粉}粒已沉降（每 \qty{15}{m}
高度的射程为 \qty{1}{km}），因此浓度
低数个数量级，\omterm{def:b2:plant-life-cycles:heterospory}{胚珠}可能要等待数日才能落上一粒\omterm{def:b2:plant-life-cycles:heterospory}{花粉}；树木自身的\omterm{def:b2:plant-life-cycles:heterospory}{花粉}
多产生自交种子，而这些种子会败育。风媒传粉需要
密集的\omterm{def:b2:plant-life-cycles:heterospory}{花粉}云，因而需要成片林分。
\textbf{18.} $10^{14}$ 粒\omterm{def:b2:plant-life-cycles:heterospory}{花粉}对应 $2\times 10^{6}$ 个\omterm{def:b2:plant-life-cycles:heterospory}{胚珠}：$5\times
10^{7}$ 粒，每个\omterm{def:b2:plant-life-cycles:heterospory}{胚珠}。
\textbf{19.} $2\times 10^{6}$ 个\omterm{def:b2:plant-life-cycles:heterospory}{胚珠} $\times 0.6\times 0.8 = 9.6\times
10^{5}$ 粒种子；\qty{4.8}{kg}；\qty{48}{MJ}。
\textbf{20.} $3\times 20/0.8 = \qty{75}{m}$：\omterm{def:b2:plant-life-cycles:heterospory}{花粉}可传播
一千米，种子只传播数个树高；基因借\omterm{def:b2:plant-life-cycles:heterospory}{花粉}传播，
植物本身则借种子传播。
\textbf{21.} $9.6\times 10^{5}\times 0.05\times 0.01 = 480$ 个成年
后代，\omterm{def:b2:plant-life-cycles:fern}{蕨类}为 57。
\textbf{22.} 松树：$4.8\times 10^{7}/480 = \qty{100}{kJ}$ 每个成年
后代；\omterm{def:b2:plant-life-cycles:fern}{蕨类}：$820/1.9 = \qty{430}{J}$ 每个成年后代，
少两百倍。
\textbf{23.} 两者均有一半为干物质，按 \qty{20}{kJ/g} 计，因此在
假设基础代谢率为每克干重 \qty{0.5}{mW}，则两者均可持续 $10^{4}\,
\text{J/g}/(5\times 10^{-4}\,\text{W/g}) = 2\times 10^{7}$ 秒，
约八个月：休眠时长并不是种子能量所
换来的。它换来的是萌发之后的过程——根与芽
在黑暗中由贮藏物质建成，早于幼苗自养数周
自养之前——并由种皮同时提供保护；孢子的
\qty{1.6e-4}{J} 仅能建成数个细胞，且须立即进行光合作用。
\textbf{24.} 种子：受精无需水；具贮备的
、受保护的休眠胚，能在干燥或季节性
环境中定居；借翅、果实和动物扩散。\omterm{def:b2:plant-life-cycles:fern}{蕨类}策略
在潮湿、遮荫、稳定的栖息地中占优，并在遥远或新
开辟的土地上定殖，那里廉价、轻巧、数量众多的繁殖体比
贮备更为重要。
\textbf{25.} \omterm{def:b2:plant-life-cycles:fern}{蕨类}：每季 $5.1\times 10^{6}$ 个孢子，对应 1.9 株成体，
每个成年后代 $2.7\times 10^{6}$ 个孢子，各 \qty{430}{J}；
松树：$9.6\times 10^{5}$ 粒种子对应 480 株成体，每
个成年后代需 \num{2000} 粒种子，耗能 \qty{100}{kJ}。
\end{solution}
